\documentclass[letterpaper,11pt]{article}

\usepackage{latexsym}
\usepackage[empty]{fullpage}
\usepackage{titlesec}
\usepackage{marvosym}
\usepackage[usenames,dvipsnames]{color}
\usepackage{verbatim}
\usepackage{enumitem}
\usepackage[hidelinks]{hyperref}
\usepackage{fancyhdr}
\usepackage[brazilian]{babel}
\usepackage{tabularx}
\input{glyphtounicode}

% Configuração de Fonte (Padrão limpo)
\usepackage[sfdefault]{FiraSans}

\pagestyle{fancy}
\fancyhf{} % Limpa cabeçalhos e rodapés
\fancyfoot{}
\renewcommand{\headrulewidth}{0pt}
\renewcommand{\footrulewidth}{0pt}

% Ajuste de Margens
\addtolength{\oddsidemargin}{-0.5in}
\addtolength{\evensidemargin}{-0.5in}
\addtolength{\textwidth}{1in}
\addtolength{\topmargin}{-0.6in}
\addtolength{\textheight}{1.2in}

\urlstyle{same}

\raggedbottom
\raggedright
\setlength{\tabcolsep}{0in}

% Formatação das Seções
\titleformat{\section}{
  \vspace{-4pt}\scshape\raggedright\large
}{}{0em}{}[\color{black}\titlerule \vspace{-5pt}]

% Garante que o PDF seja legível para sistemas de RH (ATS)
\pdfgentounicode=1

% Comandos Customizados
\newcommand{\resumeItem}[1]{
  \item\small{
    {#1 \vspace{-2pt}}
  }
}

\newcommand{\resumeSubheading}[4]{
  \vspace{-2pt}\item
    \begin{tabular*}{0.97\textwidth}[t]{l@{\extracolsep{\fill}}r}
      \textbf{#1} & #2 \\
      \textit{\small#3} & \textit{\small #4} \\
    \end{tabular*}\vspace{-7pt}
}

\newcommand{\resumeProjectHeading}[2]{
    \item
    \begin{tabular*}{0.97\textwidth}{l@{\extracolsep{\fill}}r}
      \small#1 & #2 \\
    \end{tabular*}\vspace{-7pt}
}

\newcommand{\resumeSubHeadingListStart}{\begin{itemize}[leftmargin=0.15in, label={}]}
\newcommand{\resumeSubHeadingListEnd}{\end{itemize}}
\newcommand{\resumeItemListStart}{\begin{itemize}}
\newcommand{\resumeItemListEnd}{\end{itemize}\vspace{-5pt}}

%-------------------------------------------
%%%%%%  INÍCIO DO CURRÍCULO  %%%%%%%%%%%%%%%%
\begin{document}

%----------CABEÇALHO----------
\begin{center}
    \textbf{\Huge \scshape Arthur Gonçalves Costa} \\ \vspace{5pt}
    \small Vitória, ES, Brasil $|$ (27) 98815-6932 $|$ \href{mailto:a.ctrlplay@gmail.com}{\underline{a.ctrlplay@gmail.com}} \\ \vspace{2pt}
    \href{https://www.linkedin.com/in/arthur-gon\%C3\%A7alves-costa-68153930b/}{\underline{LinkedIn}} $|$ 
    \href{https://github.com/lovssthzn}{\underline{GitHub}}
\end{center}

%----------RESUMO PROFISSIONAL----------
\section{Resumo Profissional}
\vspace{2pt}
\small{Estudante de Análise e Desenvolvimento de Sistemas (ADS) com sólida base em lógica de programação e pensamento computacional. Possuo experiência prática como responsável técnico por automações, desenvolvimento full-stack e infraestrutura, atuando de forma autônoma na estruturação de fluxos operacionais e ferramentas internas. Utilizo inteligência artificial de forma ativa e estratégica como alavancagem para explorar, compreender e implementar tecnologias, com foco constante na entrega de soluções para problemas reais de negócio.}

%----------COMPETÊNCIAS----------
\section{Competências e Habilidades}
 \begin{itemize}[leftmargin=0.15in, label={}]
    \small{\item{
     \textbf{Fundamentos:}{ Lógica de Programação, Pensamento Computacional, Resolução de Problemas, Engenharia de Prompts com IA.} \\
     \textbf{Tecnologias:}{ Ecossistema e Desenvolvimento Assistido por IA (Python, JS, React), Python, Supabase, n8n, Evolution API, Playwright, Selenium, Docker, Linux.} \\
     \textbf{Infraestrutura e Ferramentas:}{ Git, GitHub, PostgreSQL, SQL, REST API, VPS, Cloudflare, Vercel, Portainer, Web Scraping, CRM (RD Station).} \\
     \textbf{Inteligência Artificial:}{ LLMs, Retrieval-Augmented Generation (RAG), Model Context Protocol (MCP), Claude CLI, Obsidian (PKM), Processamento de Dados Não Estruturados.}
    }}
 \end{itemize}

%----------EXPERIÊNCIA----------
\section{Experiência Profissional}
  \resumeSubHeadingListStart

    \resumeSubheading
      {Assistente de Automação}{Setembro 2025 -- Março 2026}
      {CredFácil}{Vitória, ES}
      \resumeItemListStart
        \resumeItem{Atuei como o principal responsável técnico pela programação e infraestrutura da operação com total autonomia, mantendo o site institucional e desenvolvendo lógicas computacionais com suporte de IA.}
        \resumeItem{Desenvolvi scripts avançados de automação de navegador utilizando Playwright e Selenium para navegação autônoma em portais bancários (ex: Banco Mercantil), viabilizando simulações de crédito e geração de consentimentos.}
        \resumeItem{Construí fluxos automatizados via WhatsApp para triagem de leads de servidores públicos utilizando n8n e Evolution API, estruturando saídas classificadas em "enviado", "não enviado" ou "dificuldade técnica".}
        \resumeItem{Projetei e otimizei o sistema "Ponto Online" em React, garantindo a inclusão de campos de entrada manual nos botões de ação e gerenciando a conexão robusta com o banco de dados Supabase.}
        \resumeItem{Gerenciei a estabilidade e escalabilidade da operação em servidores VPS, administrando múltiplos containers em produção via Docker e Portainer e integrando APIs bancárias via REST.}
      \resumeItemListEnd

    \resumeSubheading
      {Assistente Administrativo}{Março 2025 -- Agosto 2025}
      {Facilita Brasil Oficial}{Vila Velha, ES}
      \resumeItemListStart
        \resumeItem{Atuação em rotinas administrativas e suporte à equipe comercial, utilizando diariamente o CRM RD Station para registro e acompanhamento de operações, evidenciando facilidade na adoção de softwares corporativos.}
      \resumeItemListEnd

    \resumeSubheading
      {Estagiário (Ensino Médio)}{Maio 2024 -- Janeiro 2025}
      {Prefeitura Municipal de Vitória (PMV)}{Vitória, ES}
      \resumeItemListStart
        \resumeItem{Suporte administrativo e atendimento ao público, prestando auxílio na utilização de tecnologias de comunicação e operando o sistema de prontuário eletrônico (Rede Bem Estar) e planilhas de controle.}
      \resumeItemListEnd

  \resumeSubHeadingListEnd

%----------PROJETOS----------
\section{Projetos de Destaque}
    \resumeSubHeadingListStart
      \resumeProjectHeading
          {\textbf{Gestor ITC (Instituto Todos os Cantos)} $|$ \emph{React, Vite, Tailwind, Supabase, PostgreSQL}}{2026}
          \resumeItemListStart
            \resumeItem{Desenvolvimento full-stack de um sistema web completo para a gestão administrativa do instituto, otimizando fluxos diários de trabalho.}
            \resumeItem{Implementação de módulos dedicados para controle financeiro, gestão patrimonial, acompanhamento de tarefas, agenda e administração de fornecedores.}
            \resumeItem{Estruturação de banco de dados relacional e integração de regras de negócio de acesso, permissões de usuários e relatórios operando no ecossistema Supabase.}
          \resumeItemListEnd

      \resumeProjectHeading
          {\textbf{Automação de Afiliados Shopee} $|$ \emph{Integração de APIs e Homelab}}{2026}
          \resumeItemListStart
            \resumeItem{Desenvolvimento de um sistema em Python operando 24/7 com rotinas de tempo programadas para consultar a API da Shopee, filtrando produtos de alto desempenho.}
            \resumeItem{Hospedagem da arquitetura em um servidor local (Homelab), integrando o banco de dados a fluxos do n8n para o disparo automatizado em grupos de WhatsApp.}
          \resumeItemListEnd
  

      \resumeProjectHeading
          {\textbf{Sistema RAG Local com Claude CLI e Obsidian} $|$ \emph{Arquitetura de IA}}{2026}
          \resumeItemListStart
            \resumeItem{Desenvolveu um sistema de Retrieval-Augmented Generation (RAG) local integrando Claude Code CLI e Obsidian, automatizando a ingestão de dados multimodais para consultas em linguagem natural.}
            \resumeItem{Arquitetou fluxos de extração e estruturação de conhecimento via Model Context Protocol (MCP), otimizando o contexto fornecido aos LLMs e reduzindo o tempo de pesquisa.}
          \resumeItemListEnd

      \resumeProjectHeading
          {\textbf{Homelab e Infraestrutura Pessoal} $|$ \emph{Estudo Prático}}{2026}
          \resumeItemListStart
            \resumeItem{Montagem de ambiente de servidor local, administrando serviços em containers via Docker e Portainer, e configurando túneis Cloudflare para gerenciamento de webhooks locais no n8n.}
          \resumeItemListEnd
    \resumeSubHeadingListEnd

%----------FORMAÇÃO----------
\section{Formação Acadêmica}
  \resumeSubHeadingListStart
    \resumeSubheading
      {Estácio de Sá}{Vila Velha, ES}
      {Ensino Superior em Análise e Desenvolvimento de Sistemas (ADS)}{Jun. 2025 -- Cursando}
  \resumeSubHeadingListEnd

%----------IDIOMAS----------
\section{Idiomas}
 \begin{itemize}[leftmargin=0.15in, label={}]
    \small{\item{
     \textbf{Português:}{ Nativo.} \\
     \textbf{Inglês:}{ Leitura técnica.}
    }}
 \end{itemize}

\end{document}